\section{Data and Deseasonalization}\label{sec:deseasonalization}

\subsection{Common Setup}\label{ssec:deseasonalization-common}

Each raw series is first mapped to a coordinate in which an additive model is
appropriate. Denote this transformed observable by \(Y^F_t\), where
\(F\in\{P,T,Q,W\}\) refers to price, temperature, solar, and wind. We write
\begin{equation}
  Y^F_t=\Lambda^F_t+X^F_t .
  \label{eq:deseasonalization-common-decomposition}
\end{equation}
Here \(\Lambda^F_t\) is a deterministic seasonal component and \(X^F_t\) is the
stochastic residual used in the CARMA calibration. This section specifies the
coordinate \(Y^F_t\), the seasonal structure used for \(\Lambda^F_t\), and the
diagnostic evidence left after deseasonalization. The remaining serial
dependence is then modeled stochastically.

For the physical factors \(T\), \(Q\), and \(W\), the seasonal component is a
linear regression on calendar and Fourier covariates after the relevant
transformation has been applied, with coefficients estimated by OLS. Prices use
a separate normalized seasonal construction, described below, because the
intraday price profile is more irregular.

\subsection{Electricity Prices}\label{ssec:price-deseasonalization}

We use hourly German day-ahead prices for the DE--LU bidding zone. The CARMA
calibration sample is 2023--2025 in UTC, while
\cref{fig:price-window-2021-2025} also shows 2021--2022. As seen in the figure,
the 2023--2025 window is used to exclude the exceptional price spikes associated
with the energy crisis and the war in Ukraine.

Since the spot price \(S_t\) can be negative, the stochastic price factor is
defined in shifted log coordinates,
\begin{equation}
  \widetilde S_t=\log(S_t+\delta),
  \qquad
  \delta=1000\ \text{EUR/MWh},
  \qquad
  S_t=\exp(\widetilde S_t)-\delta .
  \label{eq:price-log-shift}
\end{equation}
Thus \(Y^P_t=\widetilde S_t\). The price seasonality follows the
electricity-price deseasonalization approach discussed in \citep{Paraschiv2013}.
We estimate \(\Lambda^P_t\) with a multiplicative seasonal weight in log-price
coordinates,
\begin{equation}
  \widehat\Lambda^P_t=\bar p\,\widehat F^Y_{d(t)}\,\widehat F^D_t .
  \label{eq:price-seasonality}
\end{equation}
The daily factor \(\widehat F^Y\) is fitted from normalized daily mean
log-prices using weekday and month effects, with August split into early and
late subperiods. The intraday factor \(\widehat F^D\) is fitted from hour-of-day
profiles within 20 calendar classes, then normalized so that it only captures
shape. The estimated residual is the price factor used in the CARMA
calibration.

\begin{figure}[!htbp]
  \centering
  \includegraphics[width=0.97\textwidth]{deseasonalization/price_window_2021_2025.pdf}
  \caption{Hourly German day-ahead electricity prices from 2021 to 2025.
  The shaded region is the 2023--2025 calibration window.}
  \label{fig:price-window-2021-2025}
\end{figure}

\begin{figure}[!htbp]
  \centering
  \includegraphics[width=0.97\textwidth]{deseasonalization/price_acf_before_after_240.pdf}
  \caption{Empirical autocorrelation of the shifted log-price and of the
  deseasonalized price residual over 240 hourly lags.}
  \label{fig:price-acf-before-after}
\end{figure}

\FloatBarrier

\subsection{Temperature}\label{ssec:temperature-deseasonalization}

We use hourly 2-meter Berlin temperature over the same 2023--2025 UTC sample.
No transform is needed, so \(Y^T_t=T_t\). The seasonal regression uses a compact
Fourier basis, where \(\tau_t\) denotes the elapsed number of hours and
\(P_y=365.25\times24\).
\begin{equation}
  \Lambda^T_t
  =
  \beta_0+\beta_1\tau_t
  +\sum_{k=1}^{3}\left[
    a_k(t)\cos\left(\frac{2\pi k\tau_t}{24}\right)
    +b_k(t)\sin\left(\frac{2\pi k\tau_t}{24}\right)
  \right]
  +a_y\cos\left(\frac{2\pi\tau_t}{P_y}\right)
  +b_y\sin\left(\frac{2\pi\tau_t}{P_y}\right).
  \label{eq:temperature-seasonality}
\end{equation}
The daily-cycle coefficients \(a_k(t)\) and \(b_k(t)\) are annual Fourier
functions, so the intraday temperature profile can vary smoothly over the year.
\Cref{fig:temperature-acf-before-after} shows the residual dependence left for
the CARMA model.

\begin{figure}[!htbp]
  \centering
  \includegraphics[width=0.97\textwidth]{deseasonalization/temperature_acf_before_after_240.pdf}
  \caption{Empirical autocorrelation of observed temperature and of the
  deseasonalized temperature residual over 240 hourly lags.}
  \label{fig:temperature-acf-before-after}
\end{figure}

\FloatBarrier

\subsection{Solar Capacity Factors}\label{ssec:solar-deseasonalization}

Let \(Q_t\in[0,1]\) be the German solar capacity factor. Solar cannot be treated
as a simple additive bounded process because production is structurally zero at
night and the daylight support changes over the year. We first estimate an
empirical clear-sky envelope \(\widehat C_t=\widehat C_{d(t),h(t)}\) from the
98\% quantile of \(Q_t\) by day-of-year and hour-of-day cells, with missing-cell
completion and circular 21-day smoothing.

The stochastic solar variable is defined as a clear-sky shortfall,
\begin{equation}
\begin{aligned}
  R_t
  =
  \operatorname{clip}_{[0,1]}
  \left(1-\frac{Q_t}{\max(\widehat C_t,10^{-4})}\right),
  \qquad
  Z_t
  &=
  \operatorname{clip}_{[10^{-4},1-10^{-4}]}
  \left(\frac{R_t-10^{-4}}{0.9998}\right),
  \\
  Y^Q_t
  &=
  \log\left(\frac{Z_t}{1-Z_t}\right).
\end{aligned}
\label{eq:solar-logit-transform}
\end{equation}
A positive solar latent shock therefore means larger cloudiness or solar
shortfall, not larger physical production. Deseasonalization is performed with
\(Y^Q_t\) as the transformed coordinate. The seasonal regression combines hour
and month effects, annual Fourier terms, and hour-by-annual interactions. We
write \(H_h(t)=\mathbf 1_{\{h(t)=h\}}\), \(M_m(t)=\mathbf
1_{\{\operatorname{month}(t)=m\}}\), and \(\phi_t=2\pi d(t)/365.25\). The
regression is
\begin{equation}
  \begin{aligned}
  \Lambda^Q_t
  =
  \gamma_0
  &+
  \sum_{h=1}^{23}\gamma_h H_h(t)
  +
  \sum_{m=2}^{12}\eta_m M_m(t)
  +
  \sum_{k=1}^{3}
  \left[
    a_k\cos(k\phi_t)+b_k\sin(k\phi_t)
  \right] \\
  &+
  \sum_{h=1}^{23}H_h(t)
  \sum_{k=1}^{2}
  \left[
    u_{h,k}\cos(k\phi_t)+v_{h,k}\sin(k\phi_t)
  \right].
  \end{aligned}
  \label{eq:solar-seasonality}
\end{equation}
For physical diagnostics, the fitted latent seasonality is mapped back through
the inverse clear-sky transform.

\begin{figure}[!htbp]
  \centering
  \includegraphics[width=0.97\textwidth]{deseasonalization/solar_acf_before_after_240.pdf}
  \caption{Empirical autocorrelation of the latent solar coordinate and of the
  deseasonalized solar residual over 240 hourly lags.}
  \label{fig:solar-acf-before-after}
\end{figure}

\FloatBarrier

\subsection{Wind Capacity Factors}\label{ssec:wind-deseasonalization}

Let \(W_t\in[0,1]\) be the German wind capacity factor, aligned with the price
calendar in UTC. Wind has no structural day-night support, but it is bounded and
persistent. We therefore work in logit coordinates,
\begin{equation}
  W^\varepsilon_t=\operatorname{clip}_{[10^{-4},1-10^{-4}]}(W_t),
  \qquad
  Y^W_t=\log\left(\frac{W^\varepsilon_t}{1-W^\varepsilon_t}\right),
  \qquad
  W_t=\sigma(Y^W_t).
  \label{eq:wind-logit-transform}
\end{equation}
The transform is monotone, so a positive wind latent shock corresponds to
higher-than-seasonal physical wind production. The seasonal regression includes
hour, weekday, month, and annual Fourier effects. Physical-scale diagnostics are
obtained by mapping \(\widehat\Lambda^W_t\) back through the logistic function.

\begin{figure}[!htbp]
  \centering
  \includegraphics[width=0.97\textwidth]{deseasonalization/wind_acf_before_after_240.pdf}
  \caption{Empirical autocorrelation of the logit wind capacity factor and of
  the deseasonalized wind residual over 240 hourly lags.}
  \label{fig:wind-acf-before-after}
\end{figure}

\FloatBarrier
